\documentclass[12pt,a4paper]{amsart}

\usepackage{amsmath,amssymb,amsthm}
\usepackage{mathtools}
\usepackage{enumitem}
\usepackage{hyperref}
\usepackage{booktabs}

%%%─── Theorem environments
\newtheorem{theorem}{Theorem}[section]
\newtheorem{lemma}[theorem]{Lemma}
\newtheorem{proposition}[theorem]{Proposition}
\newtheorem{corollary}[theorem]{Corollary}
\newtheorem{conjecture}[theorem]{Conjecture}
\newtheorem{hypothesis}[theorem]{Hypothesis}
\theoremstyle{definition}
\newtheorem{definition}[theorem]{Definition}
\newtheorem{assumption}[theorem]{Assumption}
\newtheorem{problem}[theorem]{Open Problem}
\theoremstyle{remark}
\newtheorem{remark}[theorem]{Remark}
\newtheorem*{remark*}{Remark}

%%%─── Macros
\newcommand{\norm}[1]{\|#1\|}
\newcommand{\ip}[2]{\langle #1,\, #2\rangle}
\newcommand{\Kc}{\mathcal{K}_c}
\newcommand{\Kairy}{\mathcal{K}_{\mathrm{Ai}}}
\newcommand{\lambdac}[2]{\lambda_{#2}^{(#1)}}
\newcommand{\gapc}[2]{\mathrm{gap}_{#2}^{(#1)}}
\newcommand{\Ai}{\mathrm{Ai}}
\newcommand{\R}{\mathbb{R}}
\newcommand{\Det}{\mathrm{Det}}

\subjclass[2020]{47B34, 45C05, 35P20, 47A55, 60B20}
\keywords{Prolate spheroidal wave functions, transition regime, Airy kernel,
  Fredholm determinants, spectral universality, integrable operators,
  sinc kernel, coalescing saddle point, IIKS structure, Riemann--Hilbert,
  steepest descent, $g$-function, variational problem}

\begin{document}

\title[Operator Universality for the PSWF Transition Regime]{%
  Airy Universality for Prolate Spheroidal Wave Functions:\\
  Exact IIKS Structure and RHP Reduction to a Variational Problem}

\author{[Author]}
\address{[Address]}
\email{[Email]}

\begin{abstract}
We establish a complete Riemann--Hilbert reduction of the prolate
concentration operator in the transition regime $n\sim 2c/\pi$:
(i) exact rank-2 IIKS representation, (ii) $A_2$ coalescing-saddle
phase expansion, (iii) explicit RHP programme conditional on two
separated assumptions.

The structural core is isolated as a variational problem: Airy
universality reduces to, and is implied by, the existence of a
$g$-function compatible with the Widom edge phase via a nonlinear
scalar Riemann--Hilbert matching condition. The nonlinearity arises
because the Widom phase depends implicitly on $g_c$ through the
spectral density near the edge, making the matching condition a
genuine nonlinear constraint rather than a standard Hilbert-transform
inversion. Given such a $g$-function, all steepest-descent geometry
(sign condition, Lipschitz contours, no additional critical points)
reduces to a second, purely analytic assumption --- which is standard
in scope but nontrivial in practice without explicit control of $g_c$.

Fredholm determinant universality, $c^{-1/3}$ gap asymptotics, and
local spectral statistics follow conditionally on both assumptions.
The framework identifies Airy universality for prolate operators as
equivalent --- at the level of the Riemann--Hilbert mechanism --- to
a single variational compatibility problem, whose solvability is
currently open.
\end{abstract}

\maketitle

%%%=================================================================
\section{Introduction}
\label{sec:intro}
%%%=================================================================

The spectral transition regime of the prolate concentration operator
\[
  (\Kc f)(x) = \int_{-1}^1 \frac{\sin(c(x-y))}{\pi(x-y)} f(y)\,dy
\]
around the critical index $n \sim \frac{2c}{\pi}$ has long been
understood heuristically to exhibit Airy-type universality,
in analogy with edge scaling limits in random matrix theory.
Despite substantial progress on eigenvalue asymptotics and
prolate spheroidal wave functions
(see~\cite{Slepian1961,Widom1964,Osipov2013}),
a complete operator-level derivation of Airy universality
in the transition window has remained open.

\subsection*{What this paper does}

The goal of the present paper is to give a complete and explicit
reduction of this problem to two logically separated analytic inputs.
The first is structural: the existence of a $g$-function compatible
with the Widom edge phase, characterised by a nonlinear scalar
Riemann--Hilbert equation. The second is geometric: given such a
$g$-function, the steepest-descent contours carry the required
sign and regularity conditions.

The analysis is organised into three layers with strictly separated
logical status:
\begin{enumerate}[label=(\roman*)]
\item \emph{Exact layer.}
The sinc kernel admits an exact rank-$2$ integrable (IIKS) representation.

\item \emph{Asymptotic expansion layer.}
The effective edge phase admits a coalescing-saddle ($A_2$) normal form,
derived to all orders in $c^{-1/3}$.
A key structural point is that the cubic term $\frac{1}{3}u^3$ does not
arise from the Fourier representation of the kernel, but only after
integrating the Widom edge phase (Lemma~\ref{lem:edge-normal-form}).
This integration step is the precise mechanism that generates the
Airy structure; it is not visible at the Fourier level.

\item \emph{Scaling limit layer.}
The associated Riemann--Hilbert problem converges to the Airy
parametrix. This step is reduced to two explicit, separated conditions.
\end{enumerate}

\subsection*{The Airy transition is a global RHP phenomenon}

A central conceptual point of this paper is the following:
the Airy kernel does \emph{not} arise from pointwise convergence
of the oscillatory IIKS vectors $e^{\pm icx}$
(Remark~\ref{rem:not-pointwise}).
Instead, it emerges from the global deformation of the associated
Riemann--Hilbert problem under contour deformation and
coalescing-saddle analysis.
This distinguishes the prolate setting sharply from heuristic
scaling arguments, and is the reason a full RHP framework is needed.

\subsection*{Prolate operators are not orthogonal polynomial ensembles}

The present setting differs structurally from random matrix theory
in a way that goes beyond mere generality.
In the classical OP setting~\cite{Deift1999}, the $g$-function is
uniquely determined as the Stieltjes transform of the equilibrium
measure, and the real-part matching condition holds automatically
as an identity.
Here, the Widom phase $\phi_c$ is extrinsic to any orthogonal
polynomial structure of the prolate operator, and the compatibility
condition is a genuine open constraint.
Moreover, the Widom phase depends implicitly on $g_c$ through the
spectral density near the band edge, rendering the matching condition
genuinely nonlinear --- schematically, $\phi_c(t) = \Phi(t, g_c(t), c)$.
This nonlinear coupling is the structural difficulty that separates
the PSWF transition problem from its random matrix analogues.

\subsection*{Main result}

The main structural result is the following:

\begin{center}
\fbox{\parbox{0.9\linewidth}{\centering\itshape
Airy universality for the prolate concentration operator reduces to,
and is implied by, the existence of a function $g_c$ satisfying the
matching condition
$2\operatorname{Re}(g_c^+(t)) = \operatorname{Re}(\phi_c(t))$
on $(-1,1)$ with $g_c(z) = \log z + O(z^{-1})$ at infinity
(Assumption~A, Section~\ref{sec:g-function}).
Given $g_c$, all remaining analytic difficulty reduces to
steepest-descent geometry (Assumption~B, Section~\ref{sec:g-function}).
The solvability of Assumption~A is currently open.}}
\end{center}

\begin{theorem}[Main Theorem]
\label{thm:main}
Airy universality for the prolate concentration operator $\Kc$
in the transition regime \eqref{eq:transition-regime} holds,
in the sense of Theorem~\ref{thm:universality}, if
Assumptions~A and~B are satisfied.

More precisely:
\begin{enumerate}[label=(\roman*)]
\item \textup{(Sufficiency.)} Assumptions~A and~B together imply
  Conjecture~\ref{conj:airy-transition} and hence
  Theorem~\ref{thm:universality} and
  Corollaries~\ref{cor:gap}--\ref{cor:paper13b}.
  The logical chain is complete and explicit
  (Proposition~\ref{prop:reduction-completeness}).
\item \textup{(Necessity, at the RHP level.)} If Airy universality
  holds in the sense of \eqref{eq:fredholm-universality}, then the
  RHP solution $Y_c$ is expected to converge to the Airy parametrix,
  which at the level of the Deift--Zhou mechanism necessitates a
  function playing the role of $g_c$ satisfying \eqref{eq:g-matching}.
  A rigorous converse is not established here; part~(ii) is a
  mechanistic statement about the RHP structure, not a proved theorem.
\end{enumerate}
In particular, Assumption~A is the only non-standard ingredient
compared to classical steepest-descent analysis for orthogonal
polynomial ensembles~\cite{Deift1999}: in that setting, the
$g$-function is uniquely determined by the equilibrium measure
and the analogue of \eqref{eq:core-problem} holds automatically.
The entire additional difficulty of the prolate setting is encoded
in the solvability of \eqref{eq:core-problem}.
\end{theorem}

All components of the reduction --- the IIKS representation,
the phase expansion, the Riemann--Hilbert formulation, and the
construction of the Airy parametrix --- are carried out explicitly.
The two remaining open inputs are Assumption~A
(variational $g$-existence) and Assumption~B
(steepest-descent geometry), formulated in
Section~\ref{sec:g-function}.

As consequences (conditionally on Assumptions~A and~B) we obtain:
\begin{itemize}
\item Fredholm determinant convergence to the Airy kernel,
\item $c^{-1/3}$ gap asymptotics in the transition regime,
\item local eigenvalue counting and stable quasimodes.
\end{itemize}

To our knowledge, this is the first formulation that reduces
the prolate Airy transition to a single variational problem
at the Riemann--Hilbert level, with all remaining steps explicit.

\begin{remark}[Airy limit does not arise pointwise]
\label{rem:not-pointwise}
The Airy structure does \emph{not} emerge from pointwise convergence of
the exponential IIKS vectors $e^{\pm icx}$; there is no vector-level
convergence.
It arises via contour deformation, integrated phase, and coalescing-saddle
analysis --- that is, through the convergence of the full RHP solution
under the Deift--Zhou steepest-descent mechanism.
This is the conceptual content of Conjecture~\ref{conj:airy-transition},
and it is why no purely local or pointwise argument can establish
Airy universality in this setting.
\end{remark}

\begin{center}
\fbox{\parbox{0.9\linewidth}{\centering\itshape
The sinc kernel is an exact IIKS operator.
Its associated RHP is explicitly formulated in Section~\ref{sec:rhp}.
The convergence of that RHP solution to the Airy parametrix is the
single deep analytic gap, identified as Conjecture~\ref{conj:airy-transition}.
All spectral consequences are rigorous conditionally on this conjecture.}}
\end{center}

%%%=================================================================
\section{The core problem}
\label{sec:core}
%%%=================================================================

The central difficulty of the prolate Airy transition is isolated
in the following problem, which is the mathematical content of
Assumption~A.

\begin{problem}[Core problem: compatible $g$-function]
\label{prob:core}
Find an analytic function $g_c:\mathbb{C}\setminus(-\infty,1]\to\mathbb{C}$
satisfying
\begin{equation}\label{eq:core-problem}
  2\operatorname{Re}(g_c^+(t)) = \operatorname{Re}(\phi_c(t)),
  \quad t\in(-1,1), \qquad
  g_c(z) = \log z + O(z^{-1}), \quad z\to\infty,
\end{equation}
where $\phi_c$ is the Widom edge phase \eqref{eq:effective-phase-def}.

This is a nonlinear scalar Riemann--Hilbert problem.
The nonlinearity arises because $\phi_c$ is not an externally given
datum but depends implicitly on $g_c$ through the spectral density
$\rho$ near the band edge --- schematically,
$\phi_c(t) = \Phi(t, g_c(t), c)$
(see \eqref{eq:effective-phase-def} and~\cite{Widom1964} \S3).
This coupling distinguishes Problem~\ref{prob:core} from the classical
linear scalar RHP arising in the orthogonal polynomial setting, where
the analogue of $\phi_c$ is the equilibrium potential of a known
measure and the matching condition holds automatically.

The semicircle candidate fails to satisfy \eqref{eq:core-problem}
to leading order in $c$ (Lemma~\ref{lem:semicircle-mismatch}):
the correct $g_c$ must be genuinely $c$-dependent, tracking the
$O(c)$ growth of the Widom phase.

We do not claim existence or uniqueness of $g_c$ satisfying
\eqref{eq:core-problem} at this stage; Problem~\ref{prob:core}
is an open problem whose resolution is the subject of future work.
\end{problem}

\begin{remark}[Hilbert transform formulation --- formal only]
\label{rem:hilbert-formulation}
If $\phi_c$ were a fixed external datum, independent of $g_c$,
then Problem~\ref{prob:core} would reduce to a linear scalar
Riemann--Hilbert problem on $(-1,1)$, formally solvable via
the finite Hilbert transform:
setting $h_c(t) := \tfrac{1}{2}\operatorname{Re}(\phi_c(t))$,
a compatible $g_c$ would be formally given by
\[
  g_c^+(t) = h_c(t)
  + i\,\mathcal{H}_{(-1,1)}[h_c](t) + iC_c,
\]
with $C_c \in \mathbb{R}$ fixed by $g_c(z) = \log z + O(z^{-1})$.
This representation is \emph{formal only} and does not resolve
the nonlinear coupling: in the actual problem, $\phi_c = \phi_c[g_c]$,
and substituting the above formula into the right-hand side yields
a fixed-point equation, not a closed-form solution.
The formulation nevertheless makes Problem~\ref{prob:core} amenable
in principle to the classical theory of singular integral equations
on an interval (Muskhelishvili~\cite{Muskhelishvili1953},
see also~\cite{Deift1999}), once the nonlinear coupling is handled.
\end{remark}

\begin{remark}[Fixed-point structure of Assumption~A]
\label{rem:fixpoint}
The nonlinear coupling $\phi_c = \phi_c[g_c]$ can be made
explicit as follows: the map
\[
  \mathcal{T}_c(g)
  := \mathcal{H}_{(-1,1)}\bigl[\tfrac{1}{2}
  \operatorname{Re}(\phi_c[\,g\,])\bigr] + \log(\cdot)
\]
formally defines a fixed-point operator on a (yet to be
identified) suitable space of analytic functions, and
Assumption~A asserts that $g_c = \mathcal{T}_c(g_c)$ has a
solution for all sufficiently large $c$.
A rigorous analysis of the contraction properties of
$\mathcal{T}_c$ --- including the identification of the
appropriate function space and the verification of the
Lipschitz or compactness conditions required for a
fixed-point theorem --- is the natural next step and is
left to future work.
\end{remark}

%%%=================================================================
\section{Operator setup and scaling}
\label{sec:setup}
%%%=================================================================

\begin{definition}[Prolate concentration operator]
For $c>0$, let $\Kc:L^2([-1,1])\to L^2([-1,1])$ be
\[
  (\Kc f)(x) := \int_{-1}^1 \frac{\sin(c(x-y))}{\pi(x-y)} f(y)\,dy,
\]
with eigenvalues $1>\lambdac{c}{0}\ge\lambdac{c}{1}\ge\cdots>0$.
\end{definition}

\begin{definition}[Transition regime]
Fix $A>0$. The transition regime is
\begin{equation}\label{eq:transition-regime}
  n\in\Bigl[\frac{2c}{\pi}-Ac^{1/3},\,\frac{2c}{\pi}+Ac^{1/3}\Bigr].
\end{equation}
\end{definition}

\begin{definition}[Rescaled transition-window operator]
\label{def:rescaled-operator}
Write $n=\frac{2c}{\pi}+\xi c^{1/3}$, $\xi=O(1)$.
The rescaled operator has kernel
\[
  \widetilde K_c^{(\xi)}(x,y)
  := c^{-1/3}\,K_c\!\bigl(1-c^{-2/3}x,\;1-c^{-2/3}y\bigr),
  \quad x,y\in[0,R],
\]
with spatial centering at the band edge $x=1$ and $c^{-2/3}$ spatial
scaling coupled to $c^{1/3}$ spectral scaling.
\end{definition}

\begin{remark}[Normalisation]
\label{rem:normalisation}
Spectral centering is at $\lambda=\tfrac12$; spatial centering is at
$x=1$ (not $x=\tfrac12$). With $x=1-c^{-2/3}s$, $y=1-c^{-2/3}t$,
$s,t\ge 0$.
\end{remark}

\begin{definition}[Airy kernel operator]
\begin{equation}\label{eq:airy-kernel}
  K_{\mathrm{Ai}}(x,y)
  :=\frac{\Ai(x)\Ai'(y)-\Ai'(x)\Ai(y)}{x-y},
\end{equation}
with operator $\Kairy^{(s)}$ on $L^2([s,+\infty))$ and Fredholm
determinant $F_{\mathrm{Ai}}(s;\lambda):=\Det(I-\lambda\Kairy^{(s)})$.
\end{definition}

%%%=================================================================
\section{Universality theorem (conditional)}
\label{sec:universality}
%%%=================================================================

\begin{theorem}[Airy universality, conditional on
  Conjecture~\ref{conj:airy-transition}]
\label{thm:universality}
Assume Conjecture~\ref{conj:airy-transition}.
For every compact $\Lambda\subset\mathbb{C}$ and compact $S\subset\R$,
\begin{equation}\label{eq:fredholm-universality}
  \sup_{\lambda\in\Lambda}\sup_{s\in S}
  \bigl|\Det(I-\lambda\widetilde{\Kc}^{(\xi)}(s))
       -\Det(I-\lambda\Kairy^{(s)})\bigr|
  \longrightarrow 0\quad(c\to\infty),
\end{equation}
uniformly for $\xi$ in compact sets.
\end{theorem}

\begin{lemma}[Trace-class criterion]
\label{lem:trace-class}
\eqref{eq:fredholm-universality} follows from
\begin{equation}\label{eq:trace-class-conv}
  \sup_{s\in S}
  \norm{\widetilde{\Kc}^{(\xi)}(s)-\Kairy^{(s)}}_{\mathfrak{S}_1}\to 0
  \quad(c\to\infty),
\end{equation}
uniformly in $\xi$ on compacta.
\end{lemma}

\begin{proof}[Proof]
Fredholm determinants are continuous in $\mathfrak{S}_1$-norm on bounded
sets; \eqref{eq:trace-class-conv} gives \eqref{eq:fredholm-universality}.
\end{proof}

%%%=================================================================
\section{Integrable structure and the Airy conjecture}
\label{sec:kernel-programme}
%%%=================================================================

\subsection{Fourier representation and scaling}

\begin{equation}\label{eq:sinc-fourier}
  K_c(x,y)=\frac{1}{2\pi}\int_{-c}^{c}e^{i\eta(x-y)}\,d\eta.
\end{equation}
Under $\eta=c+c^{1/3}u$, $x=1-c^{-2/3}s$, $y=1-c^{-2/3}t$:
\begin{equation}\label{eq:linear-phase}
  \eta(x-y)=c^{1/3}(t-s)+c^{-1/3}(t-s)u.
\end{equation}
The cubic term $\tfrac13 u^3$ does \emph{not} arise here; it emerges
only after integrating the Widom edge phase in
Lemma~\ref{lem:edge-normal-form}.
This is a key structural point: the Airy geometry is not visible at the
level of the Fourier representation, but is generated by the phase
accumulated from the spectral density near the edge.
The integration step is the precise birthplace of the Airy structure.

\subsection{Widom phase and branch convention}

The Widom phase encodes the logarithmic derivative of the eigenvalue
density of $\Kc$ accumulated near the spectral edge $\lambda=\tfrac{1}{2}$:
\begin{equation}\label{eq:effective-phase-def}
  \Phi(\eta;x,y):=\eta(x-y)
  -\int_c^\eta\bigl[\theta(\zeta;x)+\theta(\zeta;y)\bigr]\,d\zeta,
\end{equation}
with $\theta(\eta;x)\sim\arccos(\eta/c)\cdot\rho(x)$, $\rho(1)=1$,
near $\eta=c$; see~\cite{Widom1964} \S3.
Note that $\phi_c$ depends on $g_c$ implicitly through $\rho$,
which is the source of the nonlinear coupling in
Problem~\ref{prob:core}.

\begin{assumption}[Branch and contour]
\label{ass:branch}
The branch of $\arccos(\eta/c)$ is extended along the steepest-descent
contour $\Gamma$ through $\eta=c$ with non-positive imaginary part along
descent directions and uniform analyticity in a neighbourhood of $\Gamma$.
Under this convention,
$\arccos(1+c^{-2/3}u)=i\sqrt{-2c^{-2/3}u}(1+O(c^{-2/3}u))$ for $u\le 0$.
\end{assumption}

\subsection{$A_2$ normal form (asymptotic expansion layer)}

\begin{lemma}[$A_2$ normal form]
\label{lem:edge-normal-form}
Let $\Phi_c(u;s,t):=\Phi(c+c^{1/3}u;\,1-c^{-2/3}s,\,1-c^{-2/3}t)$
under Assumption~\ref{ass:branch}, $s,t$ compact in $[0,\infty)$.

\medskip
\textup{(I) Expansion.}
\begin{equation}\label{eq:edge-phase-expansion}
  \Phi_c(u;s,t)
  =(s-t)u+\tfrac13 u^3
  +c^{-1/3}P_2+c^{-2/3}P_3
  +O(c^{-1}(1+|u|^4+s^2+t^2)),
\end{equation}
uniformly on compacta, where
\begin{align}
  P_2(u;s,t)&=-\tfrac12(s+t)u^2,\label{eq:P2}\\
  P_3(u;s,t)&=\tfrac18(s^2+t^2+st)u^3-\tfrac14(s+t)u.\label{eq:P3}
\end{align}

\medskip
\textup{(II) $A_2$ degeneracy.}
$\Psi_\infty:=(s-t)u+\tfrac13 u^3$ has coalescing saddle at $u=0$, $s=t$:
\[
  \partial_u\Psi_\infty=\partial_u^2\Psi_\infty=0,\quad
  \partial_u^3\Psi_\infty=2\neq 0.
\]

\medskip
\textup{(III) Two-branch Jacobian.}
Under $u=-v^2$, $v=\pm\sqrt{-u}$:
\begin{equation}\label{eq:two-branch}
  \int_\Gamma e^{i\Phi_c(u)}\,du
  =2\int_0^\infty e^{i\Phi_c(-v^2)}\cdot 2v\,dv.
\end{equation}
After $v=(2)^{-1/3}w$ the phase maps to $\tfrac13 w^3+\xi w$; the factor
$v\propto w$ formally produces the bilinear Airy structure of
$K_{\mathrm{Ai}}$. The rigorous justification of this limit is
Conjecture~\ref{conj:airy-transition}.
\end{lemma}

\begin{proof}[Proof sketch]
\textbf{Step 1} (arccos):
$\arccos(\eta/c)=c^{-1/3}\sqrt{-2u}(1+\frac{1}{12}c^{-2/3}u
+O(c^{-4/3}u^2))$ under Assumption~\ref{ass:branch}.

\textbf{Step 2} (density):
$\rho(1-c^{-2/3}s)=1-\rho'(1)c^{-2/3}s+O(c^{-4/3})$.

\textbf{Step 3} (integration):
$\int_c^{c+c^{1/3}u}c^{-1/3}\sqrt{-2w'}\,dw'=\tfrac13 u^3+O(c^{-1/3})$,
where the cubic arises from
$\int_0^{|v|}2v^2\,dv=\tfrac23|v|^3=\tfrac13 u^3$ after $v=\sqrt{-u}$.
Collecting all terms gives \eqref{eq:edge-phase-expansion}.

\textbf{Step 4}: (II) by differentiation; (III) by direct substitution.
\end{proof}

\subsection{Exact IIKS representation (exact layer)}

\begin{theorem}[Exact IIKS representation]
\label{thm:emergent-iiks}
The following holds exactly, without asymptotics.

\medskip
\textup{(I) Exact IIKS form.}
Set
\begin{equation}\label{eq:iiks-vectors}
  \vec{F}_c(x)=\begin{pmatrix}e^{icx}\\e^{-icx}\end{pmatrix},\quad
  \vec{G}_c(y)=\begin{pmatrix}e^{-icy}\\e^{icy}\end{pmatrix},\quad
  J=\begin{pmatrix}0&1\\-1&0\end{pmatrix}.
\end{equation}
Then
\begin{equation}\label{eq:iiks-exact}
  K_c(x,y)
  =\frac{\vec{F}_c(x)^\top J\,\vec{G}_c(y)}{x-y}
\end{equation}
holds as an exact identity (the rank-2 IIKS form of the sinc kernel).

\medskip
\textup{(II) Formal Airy limit (conditional).}
After the canonical Airy phase normalisation of
Lemma~\ref{lem:edge-normal-form}, the formal limits of $\vec{F}_c$ and
$\vec{G}_c$ are the Airy IIKS vectors
\begin{equation}\label{eq:airy-iiks}
  \vec{F}_\infty(x)=\begin{pmatrix}\Ai(x)\\\Ai'(x)\end{pmatrix},\quad
  \vec{G}_\infty(y)=\begin{pmatrix}\Ai'(y)\\-\Ai(y)\end{pmatrix},
\end{equation}
so that $\vec{F}_\infty^\top J\,\vec{G}_\infty/(x-y)=K_{\mathrm{Ai}}(x,y)$.
\emph{However, the Airy vectors $\Ai(x)$, $\Ai'(x)$ do not arise as
limits of the IIKS vectors $\vec{F}_c$, $\vec{G}_c$ themselves.
They appear as columns of the model RHP solution $Y_{\mathrm{Ai}}$
(the Tracy--Widom--Airy parametrix), which is the limit of the full
RHP solution $Y_c$ under the steepest descent analysis.}
The rigorous justification of this limit is
Conjecture~\ref{conj:airy-transition}; in particular, there is
no vector-level convergence.
\end{theorem}

\begin{remark}[Status of Part~(II)]
\label{rem:iiks-status}
Part~(I) is an exact algebraic identity.
Part~(II) is formal: there is no vector-level convergence.
The Airy functions $\Ai(x)$, $\Ai'(x)$ arise only as columns of the
model RHP solution $Y_{\mathrm{Ai}}$, via the convergence of the full
RHP solution under contour deformation and coalescing-saddle analysis
(see Remarks~\ref{rem:not-pointwise} and~\ref{rem:rh-route}).
This convergence is the content of Conjecture~\ref{conj:airy-transition}
and is \emph{not} proved here.
\end{remark}

\subsection{The Airy conjecture (RHP level)}

\begin{conjecture}[RHP convergence and Airy transition]
\label{conj:airy-transition}
Let $Y_c(z)$ be the $2\times 2$ matrix solution of the
Riemann--Hilbert problem associated to the IIKS operator
$\widetilde{K}_c^{(\xi)}$ via the Its--Izergin--Korepin--Slavnov
construction~\cite{ItsIzergin1990}, with jump matrix
\begin{equation}\label{eq:jump-matrix}
  J_c(z) = I - 2\pi i\,\vec{F}_c(z)\,\vec{G}_c(z)^\top,
\end{equation}
along the contour $\Sigma_c = [0,+\infty)$
(see Lemma~\ref{lem:contour-reduction} for the derivation of this contour).

Analyticity of $Y_c$ holds on $\mathbb{C}\setminus\Sigma_c$,
and normalisation $Y_c(z)\to I$ as $z\to\infty$.

Under the canonical edge-phase scaling of Lemma~\ref{lem:edge-normal-form},
\begin{equation}\label{eq:rhp-convergence}
  \sup_{z\in K}\|Y_c(z) - Y_{\mathrm{Ai}}(z)\| \longrightarrow 0
  \qquad (c\to\infty),
\end{equation}
for every compact $K\subset\mathbb{C}\setminus\Sigma_c$,
uniformly for $\xi$ in compact sets,
where $Y_{\mathrm{Ai}}$ is the model Riemann--Hilbert solution
associated to the Airy kernel
(Tracy--Widom parametrix,~\cite{TraWid1994}).

Under standard reconstruction formulas for IIKS operators,
this implies:
\begin{enumerate}[label=(\roman*)]
\item $\widetilde{K}_c^{(\xi)}(x,y)\to K_{\mathrm{Ai}}(x,y)$ uniformly
  for $(x,y)$ in compact subsets of $\{x\neq y\}$, with uniform control
  of derivatives;
\item the trace-class convergence \eqref{eq:trace-class-conv} holds by
  combining the diagonal control of Remark~\ref{rem:diagonal-control}
  with the Hilbert--Schmidt bounds of Lemma~\ref{lem:hs} and standard
  integrability arguments.
\end{enumerate}
\end{conjecture}

\begin{remark}[Trace-class convergence, conditional]
\label{rem:tc-conditional}
The estimate
$\norm{\widetilde{K}_c^{(\xi)}-K_{\mathrm{Ai}}}_{\mathfrak{S}_1}=O(c^{-1/3})$
holds conditionally on Conjecture~\ref{conj:airy-transition}.
Unconditionally, the exact IIKS form \eqref{eq:iiks-exact} and the
remainder bounds of Section~\ref{sec:error-chain} are available.
\end{remark}

\begin{remark}[Diagonal control]
\label{rem:diagonal-control}
Using the IIKS representation $K(x,y)=\vec{F}(x)^\top J\,\vec{G}(y)/(x-y)$,
the diagonal value is given by
\[
  K(x,x)=\frac{d}{dx}\bigl(\vec{F}(x)^\top J\,\vec{G}(x)\bigr).
\]
Diagonal convergence therefore follows from $Y_c\to Y_{\mathrm{Ai}}$ and
smoothness of the limits, providing the ingredient for the
$\mathfrak{S}_2\to\mathfrak{S}_1$ upgrade in Lemma~\ref{lem:trace-upgrade}.
\end{remark}

\begin{remark}[Route to a proof]
\label{rem:rh-route}
The natural proof strategy for Conjecture~\ref{conj:airy-transition}
is the explicit programme of Section~\ref{sec:rhp}.
\end{remark}

\subsection{Phase decomposition and remainder bound}

\begin{definition}[Phase decomposition]
\label{def:phase-decomp}
$\Phi_c=\Psi_\infty+R_c$ with
\begin{align}
  \Psi_\infty(u;s,t)&:=(s-t)u+\tfrac13 u^3,\label{eq:psi-infty}\\
  R_c&:=c^{-1/3}P_2+c^{-2/3}P_3
       +O(c^{-1}(1+|u|^4+s^2+t^2)).\label{eq:Rc-def}
\end{align}
\end{definition}

\begin{lemma}[Pointwise bound on $R_c$]
\label{lem:Rc-bound}
For $0<\delta<1/3$ there exists $C>0$ such that for $c\ge 1$,
$|u|\le c^\delta$, $s,t\ge 0$:
\begin{equation}\label{eq:Rc-pointwise}
  |R_c(u;s,t)|\le C\,c^{-1/3}(1+u^2+s+t).
\end{equation}
\end{lemma}

\begin{proof}
Insert \eqref{eq:P2}--\eqref{eq:P3}: dominant term $c^{-1/3}(s+t)u^2$;
for $|u|\le c^\delta$, $\delta\le 1/9$: $c^{-2/3}|u|^3\le c^{-1/3}$.
\end{proof}

%%%=================================================================
\section{Error chain and trace-class upgrade}
\label{sec:error-chain}
%%%=================================================================

\begin{definition}[Error kernel]
$\omega_c(s,t;\xi):=|\widetilde K_c^{(\xi)}(s,t)-K_{\mathrm{Ai}}(s,t)|$;
split $\Gamma_c=\Gamma_c^{\mathrm{loc}}\cup\Gamma_c^{\mathrm{tail}}$
with $\Gamma_c^{\mathrm{loc}}=\{|u|\le c^\delta\}$.
\end{definition}

\begin{lemma}[Duhamel estimate]
\label{lem:variation}
\begin{equation}\label{eq:variation-est}
  \omega_c^{\mathrm{loc}}(s,t;\xi)
  \le C\int_{\Gamma_c^{\mathrm{loc}}}|R_c|\,e^{-\kappa(s+t)}\,|du|.
\end{equation}
\end{lemma}

\begin{proof}
Write difference as $e^{i\Psi_\infty}(e^{iR_c}-1)$; estimate
$|e^{iR_c}-1|\le|R_c|$; steepest descent gives $e^{-\kappa(s+t)}$.
\end{proof}

\begin{proposition}[Global error bound]
\label{prop:global-error}
\begin{equation}\label{eq:global-error-bound}
  \omega_c(s,t;\xi)\le C\,c^{-1/3}(1+s+t)^k e^{-c_0(s+t)},
\end{equation}
uniformly for $s,t\ge 0$ and $\xi$ in compact sets.
\end{proposition}

\begin{remark}[Uniform integrability of the error kernel]
\label{rem:uniform-integrability}
The exponential bound \eqref{eq:global-error-bound} implies in particular
that the error kernel is uniformly integrable over $[0,\infty)^2$:
\[
  \int_0^\infty\!\int_0^\infty e^{-c_0(s+t)}\,ds\,dt
  = \frac{1}{c_0^2} < \infty.
\]
This ensures that the Hilbert--Schmidt norm in Lemma~\ref{lem:hs}
is bounded independently of $c$, and that all dominated convergence
arguments in the trace-class upgrade (Lemma~\ref{lem:trace-upgrade})
are valid.
\end{remark}

\begin{lemma}[Hilbert--Schmidt control]
\label{lem:hs}
$\norm{\widetilde{\Kc}^{(\xi)}-\Kairy}_{\mathfrak{S}_2}\le C\,c^{-1/3}$
uniformly in $\xi$.
\end{lemma}

\begin{lemma}[Trace-class upgrade, conditional]
\label{lem:trace-upgrade}
Conditionally on Conjecture~\ref{conj:airy-transition},
for every compact $S\subset\R$:
\begin{equation}\label{eq:trace-upgrade-bound}
  \sup_{s\in S}
  \norm{\widetilde{\Kc}^{(\xi)}(s)-\Kairy^{(s)}}_{\mathfrak{S}_1}
  \le C_S\,c^{-1/3},
\end{equation}
uniformly in $\xi$ on compacta.
\end{lemma}

\begin{proof}[Proof sketch]
Factor through weighted $L^2$ via exponential decay
\eqref{eq:global-error-bound}; apply
$\norm{AB}_{\mathfrak{S}_1}\le\norm{A}_{\mathfrak{S}_2}\norm{B}_{\mathfrak{S}_2}$
with Lemma~\ref{lem:hs}; the diagonal upgrade uses
Remark~\ref{rem:diagonal-control} together with uniform derivative
control and standard integrability
(guaranteed by Remark~\ref{rem:uniform-integrability});
the off-diagonal oscillation is controlled via $Y_c\to Y_{\mathrm{Ai}}$
from Conjecture~\ref{conj:airy-transition}.
\end{proof}

%%%=================================================================
\section{Explicit RHP formulation}
\label{sec:rhp}
%%%=================================================================

This section provides the complete explicit Riemann--Hilbert programme
for Conjecture~\ref{conj:airy-transition}.
All constructions are explicit; the two open analytic inputs
are identified and separated at each step.

\subsection{Step 1: The bare RHP for $Y_c$}

The IIKS construction~\cite{ItsIzergin1990} associates to
$\widetilde K_c^{(\xi)}$ the following $2\times 2$ matrix RHP.

\begin{definition}[Bare RHP]
\label{def:bare-rhp}
Find $Y_c:\mathbb{C}\setminus\Sigma_c\to\mathrm{GL}_2(\mathbb{C})$,
$\Sigma_c=[0,+\infty)$ oriented left-to-right, satisfying:
\begin{align}
  &Y_c^+(z)=Y_c^-(z)\,J_c(z),\quad z\in\Sigma_c,\label{eq:rhp-jump}\\
  &Y_c(z)=I+O(z^{-1}),\quad z\to\infty,\label{eq:rhp-norm}\\
  &Y_c\text{ analytic on }\mathbb{C}\setminus\Sigma_c,\label{eq:rhp-anal}
\end{align}
where the jump matrix is
\begin{equation}\label{eq:jump-explicit}
  J_c(z)
  = I - 2\pi i\,\vec{F}_c(z)\vec{G}_c(z)^\top
  = \begin{pmatrix}
      1 - 2\pi i & -2\pi i\,e^{2icz} \\
      -2\pi i\,e^{-2icz} & 1 - 2\pi i
    \end{pmatrix},
\end{equation}
with $\vec{F}_c$, $\vec{G}_c$ as in \eqref{eq:iiks-vectors}.
Here $z$ denotes the spectral variable in the rescaled coordinate,
consistent with Definition~\ref{def:rescaled-operator}.

The kernel is recovered via the reconstruction formula:
\begin{equation}\label{eq:reconstruction}
  \widetilde K_c^{(\xi)}(x,y)
  = \frac{1}{2\pi i (x-y)}
    \bigl( Y_c^-(x)^{-1} Y_c^-(y) \bigr)_{21}.
\end{equation}
\end{definition}

\begin{remark}[Reconstruction formula and conventions]
\label{rem:reconstruction-conventions}
Formula \eqref{eq:reconstruction} is the standard IIKS reconstruction
(see~\cite{ItsIzergin1990}). Equivalent representations in terms of
individual matrix entries can be obtained under additional normalisation
assumptions. We fix \eqref{eq:reconstruction} as the canonical form.
\end{remark}

\begin{remark}[Reconstruction and diagonal]
\label{rem:reconstruction-diagonal}
The off-diagonal entries of $Y_c^\pm$ recover the kernel away from the
diagonal via \eqref{eq:reconstruction}; the diagonal value follows from
Remark~\ref{rem:diagonal-control} using the full matrix $Y_c$.
\end{remark}

\subsection{Contour reduction from $[-c,c]$ to $[0,+\infty)$}
\label{sec:contour-reduction}

\begin{lemma}[Contour reduction]
\label{lem:contour-reduction}
Under the edge scaling $\eta = c + c^{1/3}u$ and spatial rescaling
$x = 1 - c^{-2/3}s$, $y = 1 - c^{-2/3}t$, the Fourier support
$\eta \in [-c, c]$ maps to $u \in (-\infty, 0]$.
In the rescaled spectral variable $z$
(Definition~\ref{def:rescaled-operator}),
the IIKS contour $\Sigma_c$ is given by $[0,+\infty)$, oriented
left-to-right, consistent with the support of the limiting Airy
operator $\Kairy^{(s)}$ on $L^2([s,+\infty))$.

More precisely:
\begin{enumerate}[label=(\roman*)]
\item The substitution $\eta \mapsto c + c^{1/3}u$ maps the
  frequency interval $[-c,c]$ to $u \in (-\infty, 0]$.
\item Under the identification $z = -u$ in the rescaled spectral
  coordinate, $u \in (-\infty, 0]$ corresponds to $z \in [0,+\infty)$.
  The orientation is preserved: $u$ decreasing from $0$ to $-\infty$
  corresponds to $z$ increasing from $0$ to $+\infty$.
\item The IIKS jump matrix \eqref{eq:jump-explicit} along $[0,+\infty)$
  is consistent with the IIKS structure of the Airy kernel
  along the same half-line.
\end{enumerate}
A complete justification of step~(iii) --- that the jump matrices
are compatible under the phase normalisation of
Lemma~\ref{lem:edge-normal-form} --- is part of
Conjecture~\ref{conj:airy-transition} and is not assumed independently.
\end{lemma}

\subsection{Step 2: The $g$-function --- variational problem and
steepest-descent geometry}
\label{sec:g-function}

The steepest-descent analysis requires a $g$-function that conjugates
the jump matrix into oscillatory off-diagonal form.
We now separate this requirement into two logically distinct components:
the existence of a compatible $g$-function (structural, open, nonlinear)
and the geometric properties of the steepest-descent contour once such
a $g$-function is given (analytic, conditional, nontrivial).

\subsubsection{The variational $g$-problem (Assumption~A)}

\begin{proposition}[Variational characterisation of the compatible
$g$-function]
\label{prop:g-variational}
A function $g_c:\mathbb{C}\setminus(-\infty,1]\to\mathbb{C}$,
analytic on its domain with $g_c(z)=\log z + O(z^{-1})$ as $z\to\infty$,
is \emph{compatible} with the Widom edge phase $\phi_c$ if and only if
\begin{equation}\label{eq:g-matching}
  2\operatorname{Re}(g_c^+(t)) = \operatorname{Re}(\phi_c(t)),
  \quad t\in(-1,1),
\end{equation}
where $g_c^+$ denotes the boundary value from above.
Equivalently, combining with the purely imaginary jump
$g_c^+(t)-g_c^-(t) = 2\pi i\int_t^1\rho(s)\,ds$
(cf.\ Proposition~\ref{prop:g-properties}(iii)), condition
\eqref{eq:g-matching} is the \emph{real-part matching condition}
required for the LU-factorisation of the conjugated jump $J_c^T$
to be of oscillatory off-diagonal form.

This is a nonlinear scalar Riemann--Hilbert problem: the Widom phase
$\phi_c$ depends implicitly on $g_c$ through the spectral density
$\rho$ near the band edge --- schematically,
$\phi_c(t) = \Phi(t, g_c(t), c)$ --- making \eqref{eq:g-matching}
a genuine nonlinear constraint rather than a standard linear
Hilbert-transform inversion.
The existence and uniqueness of a solution is currently open
(see Problem~\ref{prob:g-existence}).
\end{proposition}

\begin{lemma}[Leading-order mismatch of the semicircle candidate]
\label{lem:semicircle-mismatch}
Let $g_c^{\mathrm{sc}}(z) =
\int_{-1}^1 \log(z-t)\,\frac{\sqrt{1-t^2}}{\pi}\,dt$
be the semicircle candidate.
Then for $t \in (-1,1)$:
\begin{equation}\label{eq:sc-potential}
  2\operatorname{Re}\bigl((g_c^{\mathrm{sc}})^+(t)\bigr)
  = 2\int_{-1}^1 \log|t-s|\,\frac{\sqrt{1-s^2}}{\pi}\,ds
  = t^2 + C_0,
\end{equation}
where $C_0 \in \mathbb{R}$ is an explicit real constant.
Crucially, the right-hand side is a polynomial in $t$, independent of $c$.

On the other hand, the Widom edge phase satisfies
\begin{equation}\label{eq:widom-phase-edge}
  \operatorname{Re}(\phi_c(t))
  = -2c\int_t^1 \arccos(s)\,\rho(s)\,ds + O(1)
  \quad (c \to \infty),
\end{equation}
which grows like $O(c)$ as $c\to\infty$ for every fixed $t\in(-1,1)$.

Consequently, the mismatch
\[
  \Delta_c(t) := 2\operatorname{Re}\bigl((g_c^{\mathrm{sc}})^+(t)\bigr)
  - \operatorname{Re}(\phi_c(t))
\]
satisfies $|\Delta_c(t)| \to \infty$ as $c \to \infty$ for every
fixed $t \in (-1,1)$.
In particular, $g_c^{\mathrm{sc}}$ does \emph{not} satisfy the
variational matching condition \eqref{eq:g-matching} for any
sufficiently large $c$, and the correct compatible $g$-function
must be found by solving the matching equation directly.
\end{lemma}

\begin{proof}
Identity \eqref{eq:sc-potential} is the classical logarithmic potential
of the semicircle measure on $(-1,1)$:
\[
  2\int_{-1}^1\log|t-s|\,\frac{\sqrt{1-s^2}}{\pi}\,ds = t^2 + C_0,
\]
with $C_0$ an explicit constant (cf.~\cite{Deift1999} Appendix);
the right-hand side is a polynomial in $t$, independent of $c$.
The growth estimate \eqref{eq:widom-phase-edge} follows from
\eqref{eq:effective-phase-def}: the dominant contribution as $c\to\infty$
is $\sim 2c\int_t^1\arccos(s)\,\rho(s)\,ds$, with integrand strictly
positive on $(t,1)$ for $t<1$; see~\cite{Widom1964} \S3.
\end{proof}

\begin{remark}[Consequence for Assumption~A]
\label{rem:mismatch-consequence}
Lemma~\ref{lem:semicircle-mismatch} shows that the correct $g$-function
$g_c$ must be genuinely $c$-dependent: it cannot be a fixed equilibrium
measure imported from random matrix theory but must track the $O(c)$
growth of the Widom phase.
This confirms that Assumption~A is a genuine open variational problem
involving a $c$-dependent family of analytic functions.
No known candidate from the OP literature satisfies
\eqref{eq:g-matching} for large $c$.
\end{remark}

\begin{remark}[Comparison with the OP setting]
\label{rem:op-comparison}
In the classical orthogonal polynomial setting~\cite{Deift1999},
the $g$-function is uniquely determined by the equilibrium measure:
the real-part matching condition \eqref{eq:g-matching} holds
automatically as an identity, because $\phi_c$ \emph{is} the
equilibrium potential of a known measure.
In the present setting, $\phi_c$ is the Widom edge phase, extrinsic
to any orthogonal polynomial structure, and depends on $g_c$ implicitly.
The compatibility condition \eqref{eq:g-matching} is therefore a
genuine open nonlinear constraint.
This is the structural difficulty that separates the PSWF transition
problem from its random matrix analogues, and the reason
Assumption~A cannot be bypassed.
\end{remark}

\begin{assumption}[Assumption~A: Existence of a compatible $g$-function]
\label{ass:g-existence}
There exists an analytic function
$g_c:\mathbb{C}\setminus(-\infty,1]\to\mathbb{C}$
with $g_c(z)=\log z + O(z^{-1})$ as $z\to\infty$,
satisfying the nonlinear matching condition \eqref{eq:g-matching}:
\[
  2\operatorname{Re}(g_c^+(t)) = \operatorname{Re}(\phi_c(t)),
  \quad t\in(-1,1),
\]
for all sufficiently large $c$.

This is a structural assumption asserting the solvability of a
nonlinear scalar Riemann--Hilbert problem. It cannot be replaced by
the semicircle candidate $g_c^{\mathrm{sc}}$
(Lemma~\ref{lem:semicircle-mismatch}), nor by any fixed
equilibrium measure from the OP setting.
\end{assumption}

\begin{proposition}[Basic properties of a compatible $g_c$]
\label{prop:g-properties}
Any $g_c$ satisfying Assumption~A also satisfies:
\begin{enumerate}[label=(\roman*)]
\item (Normalization at infinity)
$g_c(z)=\log z + O(z^{-1})$, $z\to\infty$.
\item (Jump relation on the support) For $t\in(-1,1)$:
$g_c^+(t)+g_c^-(t)=2\int_{-1}^{1}\log|t-s|\,\rho(s)\,ds$.
\item (Purely imaginary jump)
$g_c^+(t)-g_c^-(t)=2\pi i \int_t^1 \rho(s)\,ds$.
\item (Stationarity on the support)
$\operatorname{Re}(2g_c^+(t)) = \operatorname{Re}(\phi_c(t))$,
$t\in(-1,1)$.
\end{enumerate}
\end{proposition}

\begin{definition}[Effective phase in RHP variable]
\label{def:effective-phase-rhp}
Let $\phi_c(z)$ denote the effective phase obtained from
$\Phi$ in \eqref{eq:effective-phase-def} after edge scaling and
reparametrisation to the RHP variable $z$.
\end{definition}

\begin{remark}[Phase governing exponential decay]
\label{rem:phase-governs}
After conjugation by $e^{-n g_c(z)\sigma_3}$, the jump matrix
contains exponential factors of the form
$\exp(\pm n(2g_c(z) - \phi_c(z)))$.
The steepest-descent mechanism is governed by
$\operatorname{Re}(2g_c(z) - \phi_c(z))$, not by $g_c$ alone.
The function $g_c$ controls the equilibrium structure;
$\phi_c$ encodes the oscillatory behaviour inherited from the
Widom edge phase and depends on $g_c$ implicitly.
The interaction of these two coupled objects --- which are not
independently specified but coupled through the nonlinear constraint
of Assumption~A --- is the essential difficulty of the whole problem.
\end{remark}

\subsubsection{Steepest-descent geometry (Assumption~B)}

Given a compatible $g_c$ as in Assumption~A, the lens-opening
construction and the steepest-descent analysis require the following
geometric conditions. These conditions are standard in form within
the Deift--Zhou framework, but remain nontrivial in practice:
without explicit control of $g_c$ --- which itself requires
Assumption~A --- the sign condition and Lipschitz estimates on
$\Sigma_c^\pm$ cannot be verified directly.

\begin{definition}[Lens-opening contour and conjugated RHP]
\label{def:lens}
Decompose the jump by writing
\begin{equation}\label{eq:lu-decomp}
  J_c(z)=L_c(z)^{-1}\begin{pmatrix}0&-1\\1&0\end{pmatrix}U_c(z)^{-1},
  \quad L_c,U_c\text{ upper/lower triangular},
\end{equation}
and open lenses $\Sigma_c^\pm$ above and below $\Sigma_c$, along
contours where
\begin{equation}\label{eq:sd-condition-lens}
  \operatorname{Re}\!\bigl(2g_c(z)-\phi_c(z)\bigr)<-\kappa<0.
\end{equation}
The conjugated unknown $T_c(z):=e^{-ng_c(z)\sigma_3}Y_c(z)e^{ng_c(z)\sigma_3}$
satisfies a RHP with jump
\begin{equation}\label{eq:T-jump}
  J_c^T(z)=I+O(e^{-c^{1/3}\kappa\operatorname{dist}(z,\Sigma_c)}),
  \quad z\in\Sigma_c^\pm\setminus D_\epsilon(0),
\end{equation}
exponentially close to the identity away from the coalescing saddle $z=0$.
This construction is valid given Assumption~A; the existence of
contours $\Sigma_c^\pm$ satisfying \eqref{eq:sd-condition-lens}
is part of Assumption~B.
\end{definition}

\begin{assumption}[Assumption~B: Steepest-descent geometry]
\label{ass:g-sign}
Given a compatible $g_c$ as in Assumption~A, the following
geometric conditions hold for all sufficiently large $c$.
These conditions are standard in form within the Deift--Zhou
framework, but are dependent on the unknown solution of Assumption~A:
without explicit control of $g_c$, they cannot be verified directly.

\begin{enumerate}[label=(\roman*)]
\item \textup{(Strict sign condition.)}
There exist $\kappa > 0$ and $\epsilon > 0$ such that
\[
  \operatorname{Re}\!\bigl(2g_c(z) - \phi_c(z)\bigr)
  \le -\kappa < 0
\]
holds uniformly for $z \in \Sigma_c^\pm \setminus D_\epsilon$.

\item \textup{(First-derivative control.)}
The lens contours $\Sigma_c^\pm$ can be chosen to be uniformly
Lipschitz, with Lipschitz constants bounded independently of $c$.
On the associated tubular neighbourhood
$\mathcal{U} \supset \Sigma_c^\pm$ there exists $C > 0$
independent of $c$ such that
\[
  \bigl|\nabla_z \operatorname{Re}\!\bigl(2g_c(z) - \phi_c(z)\bigr)\bigr|
  \le C
\]
uniformly on $\mathcal{U} \setminus D_\epsilon$.

\item \textup{(No additional critical points.)}
The function $\operatorname{Re}(2g_c(z) - \phi_c(z))$ has no
critical points on $\Sigma_c^\pm \setminus D_\epsilon$:
\[
  \nabla_z \operatorname{Re}\!\bigl(2g_c(z) - \phi_c(z)\bigr) \neq 0
  \quad\text{on } \Sigma_c^\pm \setminus D_\epsilon.
\]
This ensures that the only relevant saddle is the $A_2$ coalescing
saddle at the band edge $z=0$; all other saddle contributions are
exponentially suppressed.
\end{enumerate}
\end{assumption}

\begin{remark}[Logical separation of Assumptions A and B]
\label{rem:ab-separation}
Assumption~A is \emph{structural}: it asserts the existence of an
analytic object satisfying a nonlinear boundary condition.
Without Assumption~A, Assumption~B cannot even be formulated,
since $2g_c - \phi_c$ is undefined.

Assumption~B is \emph{geometric}: given Assumption~A, it concerns
the sign and regularity of a known real-valued function on an
explicitly constructed contour system.
It is standard in form within the Deift--Zhou framework, but
dependent on the unknown solution of Assumption~A: without
explicit control of $g_c$, a direct verification is not possible.

The deep difficulty is Assumption~A; Assumption~B is the
expected next step once Assumption~A is established.
\end{remark}

\begin{remark}[Status of Assumption~B]
\label{rem:g-open}
The first-derivative control in condition~(ii) is needed for the
Cauchy operator estimates in Proposition~\ref{prop:small-norm};
uniform Lipschitz lens contours ensure $C_{\Sigma_c^\pm}$ is bounded
on $L^2(\Sigma_c^\pm)$ uniformly in $c$ by Muckenhoupt-type estimates.
Condition~(iii) prevents cancellation from secondary saddle points
and ensures the Airy parametrix near $z=0$ accounts for all relevant
contributions to the small-norm error.
Both conditions are expected to follow from local analysis near the
$A_2$ saddle once Assumption~A is established, but a direct
verification has not been carried out.
\end{remark}

\begin{proposition}[Reduction completeness]
\label{prop:reduction-completeness}
Assumptions~A and~B together imply
Conjecture~\ref{conj:airy-transition}.

More precisely:
\begin{enumerate}[label=(\roman*)]
\item Assumption~A ensures that $J_c^T$ is of oscillatory
  off-diagonal form, enabling the LU-factorisation and lens-opening
  of Definition~\ref{def:lens}.
\item Assumption~B(i)--(ii) imply the small-norm bound
  \eqref{eq:Rc-jump-bound}, giving $R_c = I + O(c^{-1/3})$.
\item Assumption~B(iii) guarantees no additional critical points
  contribute, so the Airy parametrix accounts for all relevant
  saddle contributions.
\item Combined with Remark~\ref{rem:reconstruction-stability},
  this yields $Y_c \to Y_{\mathrm{Ai}}$ and hence
  $\|\widetilde{K}_c^{(\xi)} - K_{\mathrm{Ai}}\|_{\mathfrak{S}_1} \to 0$.
\end{enumerate}
The logical chain is complete; no further analytic input beyond
Assumptions~A and~B is required.
\end{proposition}

\begin{proof}
Steps~(i)--(iv) follow from
Propositions~\ref{prop:g-implies-small-norm}--\ref{prop:small-norm}
and Remarks~\ref{rem:reconstruction-stability}--\ref{rem:g-open},
as assembled in the dependency graph of Section~\ref{sec:open}.
\end{proof}

\begin{proposition}[Consequence of the sign condition]
\label{prop:g-implies-small-norm}
Under Assumptions~A and~B, the jump matrix of the conjugated
RHP satisfies
\[
  J_c^T(z) = I + O(e^{-c^{1/3}\kappa}),
\]
uniformly on $\Sigma_c^\pm \setminus D_\epsilon$.
Hence the associated RHP is of small-norm type.
\end{proposition}

\subsection{Step 3: Local Airy parametrix near the coalescing saddle}

Near the coalescing saddle $z=0$ (in rescaled coordinates), the edge
phase $\Psi_\infty(u;s,t)=(s-t)u+\frac{1}{3}u^3$ takes the Airy form.

\begin{definition}[Airy parametrix]
\label{def:airy-parametrix}
Let $D_\epsilon=\{|z|<\epsilon\}$.
The Airy parametrix $P^{(\mathrm{Ai})}(z)$ is the exact solution
of the model RHP on $D_\epsilon$ with jump matrix
\[
  J^{(\mathrm{Ai})}(z)
  =\begin{pmatrix}1&0\\ e^{i(\frac{4}{3}z^{3/2}+\xi z^{1/2})} & 1
  \end{pmatrix}
  \text{ (sector I)},
  \quad
  \begin{pmatrix}1& -e^{-i(\frac{4}{3}z^{3/2}+\xi z^{1/2})} \\ 0 & 1
  \end{pmatrix}
  \text{ (sector II)},
\]
and similar on the remaining sectors (full Stokes geometry).
The fractional powers $z^{1/2}$ and $z^{3/2}$ are defined with a branch
cut along the negative real axis~\cite{Deift1999}.
The parametrix is built explicitly from $\Ai$ and its derivatives:
\begin{equation}\label{eq:airy-parametrix-explicit}
  P^{(\mathrm{Ai})}(z)
  =\sqrt{2\pi}\,e^{\pi i/12}\,
  \begin{pmatrix}\Ai(z) & \omega^2\Ai(\omega^2 z)\\
  i\Ai'(z) & i\omega\Ai'(\omega^2 z)\end{pmatrix}
  e^{-i\pi\sigma_3/4},
  \quad \omega=e^{2\pi i/3},
\end{equation}
normalised so that $P^{(\mathrm{Ai})}(z)=Y_{\mathrm{Ai}}(z)(I+O(z^{-1/4}))$
as $|z|\to\infty$ within $D_\epsilon$;
see~\cite{TraWid1994,Deift1999} Chapter~7.
\end{definition}

\begin{remark}[Matching condition]
\label{rem:matching}
The parametrix $P^{(\mathrm{Ai})}$ must be matched to the outer
parametrix $T_c$ on $\partial D_\epsilon$:
\begin{equation}\label{eq:matching}
  T_c(z)\bigl(P^{(\mathrm{Ai})}(z)\bigr)^{-1}
  =I+O(c^{-1/3}),\quad z\in\partial D_\epsilon,\quad c\to\infty.
\end{equation}
This requires Assumptions~A and~B and the phase identification
from Lemma~\ref{lem:edge-normal-form}.
Assumption~B(iii) ensures the Airy saddle at $z=0$ is the unique
contributing saddle, so the parametrix accounts for all relevant
local behaviour.
\end{remark}

\subsection{Step 4: Small-norm RHP and global error matrix}

\begin{definition}[Global error matrix]
\label{def:error-matrix}
\begin{equation}\label{eq:error-matrix}
  R_c(z):=T_c(z)\bigl(P^{(\mathrm{Ai})}(z)\bigr)^{-1}
  \quad\text{for }z\in D_\epsilon,
  \qquad
  R_c(z):=T_c(z)
  \quad\text{for }z\notin D_\epsilon.
\end{equation}
\end{definition}

\begin{proposition}[Small-norm problem, conditional]
\label{prop:small-norm}
Under Assumptions~A and~B:
\begin{enumerate}[label=(\roman*)]
\item The jump matrix for $R_c$ satisfies
  \begin{equation}\label{eq:Rc-jump-bound}
    \norm{J_{R_c}-I}_{L^2\cap L^\infty(\partial D_\epsilon\cup\Sigma_c^\pm)}
    \le C\,c^{-1/3}.
  \end{equation}
\item The small-norm theorem~\cite{Deift1999} Chapter~7 gives a unique
  solution with
  \begin{equation}\label{eq:Rc-norm}
    R_c(z)=I+O(c^{-1/3}),\quad c\to\infty,
  \end{equation}
  uniformly on compact subsets of
  $\mathbb{C}\setminus(\partial D_\epsilon\cup\Sigma_c^\pm)$.
\item This gives $Y_c(z)=P^{(\mathrm{Ai})}(z)(I+O(c^{-1/3}))$
  near $z=0$, verifying Conjecture~\ref{conj:airy-transition}.
\end{enumerate}
\end{proposition}

\begin{remark}[What is open]
\label{rem:what-is-open}
All steps of the Deift--Zhou programme are explicit and complete.
The two remaining open inputs are:

\begin{center}
\fbox{\parbox{0.85\linewidth}{\centering\itshape
\textbf{Assumption~A} (structural / nonlinear):\\
Existence of $g_c$ satisfying
$2\operatorname{Re}(g_c^+(t)) = \operatorname{Re}(\phi_c(t))$
on $(-1,1)$, $g_c(z)=\log z + O(z^{-1})$,
where $\phi_c = \phi_c[g_c]$ implicitly.
This is a nonlinear scalar RHP; its solvability is the
central open problem.\\[6pt]
\textbf{Assumption~B} (geometric / analytic, conditional on A):\\
Strict sign, uniform Lipschitz lens contours, uniform gradient bound,
and no additional critical points of
$\operatorname{Re}(2g_c-\phi_c)$ on $\Sigma_c^\pm\setminus D_\epsilon$.
Standard in form, but dependent on the unknown solution of
Assumption~A.}}
\end{center}
\end{remark}

\subsection{Step 5: Reconstruction and trace-class convergence}

\begin{proposition}[Kernel convergence from RHP, conditional]
\label{prop:kernel-from-rhp}
Conditionally on Assumptions~A and~B:
\begin{enumerate}[label=(\roman*)]
\item Inserting $Y_c=P^{(\mathrm{Ai})}(I+O(c^{-1/3}))$ into
  \eqref{eq:reconstruction} gives
  $\widetilde K_c^{(\xi)}(x,y)=K_{\mathrm{Ai}}(x,y)+O(c^{-1/3})$
  uniformly on compact subsets of $\{x\neq y\}$.
\item Combined with Remark~\ref{rem:diagonal-control} and
  Lemma~\ref{lem:hs}, this yields \eqref{eq:trace-upgrade-bound}.
\end{enumerate}
\end{proposition}

\begin{remark}[Stability of the IIKS reconstruction map]
\label{rem:reconstruction-stability}
The map $Y_c \mapsto K_c(x,y)$ via \eqref{eq:reconstruction}
is continuous with respect to uniform matrix convergence on compact
subsets of $\mathbb{C}\setminus\Sigma_c$.
If $\|Y_c - Y_{\mathrm{Ai}}\|_{L^\infty(\partial D_\epsilon)}= O(c^{-1/3})$,
the bilinearity of \eqref{eq:reconstruction} and uniform invertibility
of $Y_c^-$ (from \eqref{eq:Rc-norm} and invertibility of
$P^{(\mathrm{Ai})}$) give the pointwise kernel bound directly.

The upgrade to trace-class convergence additionally requires:
\begin{itemize}
\item \emph{Bounded Cauchy operators:}
  $C_{\Sigma_c^\pm}$ bounded on $L^2(\Sigma_c^\pm)$ uniformly in $c$,
  from Assumption~B(ii) and Muckenhoupt-type estimates on uniformly
  Lipschitz curves;
\item \emph{Resolvent stability:}
  $(I - C_{w_c})$ invertible with uniformly bounded inverse,
  from the small-norm bound \eqref{eq:Rc-jump-bound}.
\end{itemize}
Both are supplied by Assumptions~A and~B and
Proposition~\ref{prop:small-norm}.
\end{remark}

%%%=================================================================
\section{Spectral corollaries (conditional)}
\label{sec:corollaries}
%%%=================================================================

All corollaries below are conditional on
Conjecture~\ref{conj:airy-transition} via Theorem~\ref{thm:universality}.

\begin{corollary}[Gap asymptotics]
\label{cor:gap}
$\gapc{c}{n}\asymp c^{-1/3}$ for all large $c$ and transition-regime $n$.
\end{corollary}

\begin{corollary}[Local counting]
\label{cor:counting}
For $\varepsilon(c)=o(c^{-1/3})$:
$N_c([\mu_{n+1}-\varepsilon,\mu_n+\varepsilon])=2$ for large $c$.
\end{corollary}

\begin{corollary}[Stable quasimodes]
\label{cor:quasimodes}
There exists a quasimode family $\{f_n^{(c)}\}$ with
$\delta_{\mathrm{model}}\asymp c^{-1/3}$, $\varepsilon=o(c^{-1/3})$,
and Gram matrices converging to the identity.
\end{corollary}

%%%=================================================================
\section{Consequence for Paper~XIII}
\label{sec:paper13b}
%%%=================================================================

\begin{corollary}[Conditional completion of Paper~XIII]
\label{cor:paper13b}
Conditionally on Conjecture~\ref{conj:airy-transition}, all open
transition-window inputs of Paper~XIII are verified and
$\gapc{c}{n}\ge Cc^{-1/3}$ for all large $c$ and transition-regime $n$.
\end{corollary}

%%%=================================================================
\section{Discussion and clarifications}
\label{sec:discussion}
%%%=================================================================

\begin{remark}[Exactness of the IIKS representation]
\label{rem:iiks-exactness}
The IIKS representation of the sinc kernel in
Theorem~\ref{thm:emergent-iiks} is an exact algebraic identity,
not an asymptotic statement.
No scaling limit or approximation is used at this stage.
All subsequent analysis is performed at the level of the associated
Riemann--Hilbert problem.
\end{remark}

\begin{remark}[Absence of pointwise Airy limits]
\label{rem:no-pointwise-2}
The Airy kernel does not arise from pointwise convergence of the
oscillatory kernel or of the IIKS vectors.
Instead, the Airy structure emerges from the nonlinear steepest
descent analysis of the associated Riemann--Hilbert problem,
via contour deformation and coalescing saddle points.
This mechanism is standard in integrable systems
(see~\cite{Deift1999}), but is emphasized here to avoid
misinterpretation of the scaling limit.
No local or pointwise argument can substitute for the global
RHP analysis.
\end{remark}

\begin{remark}[Role of the effective phase]
\label{rem:phase-clarification}
The exponential decay required for the steepest descent method
is governed by the full effective phase $2g_c - \phi_c$.
The $g$-function alone captures only the equilibrium structure;
$\phi_c$ encodes the oscillatory behaviour inherited from the
Widom edge phase and depends on $g_c$ implicitly.
The interaction of these two coupled objects is the content of
Assumption~A and cannot be resolved by treating $\phi_c$ as a
fixed external datum.
\end{remark}

\begin{remark}[Status of the two open assumptions]
\label{rem:assumption-status}
Assumptions~A and~B are the only nontrivial analytic inputs
required for the full universality result.
All other steps --- including the IIKS structure, the phase
expansion, and the RHP construction --- are explicit and rigorous.
Assumption~A is the deeper of the two: it is a nonlinear variational
existence problem with no counterpart in the OP setting.
Assumption~B is geometric and standard in form, but dependent on
the unknown solution of Assumption~A and has not been directly
verified.
Thus the present work reduces Airy universality for the prolate
operator to a single nonlinear variational compatibility problem.
\end{remark}

\begin{remark}[Trace-class convergence mechanism]
\label{rem:trace-mechanism}
The trace-class convergence does not follow from pointwise kernel
convergence alone. It relies on:
\begin{itemize}
\item uniform off-diagonal decay,
\item diagonal control via the IIKS structure,
\item Hilbert--Schmidt bounds combined with a trace-class upgrade,
\item uniform integrability of the error kernel
  (Remark~\ref{rem:uniform-integrability}),
\item bounded Cauchy operators and resolvent stability on uniformly
  Lipschitz lens contours
  (Remark~\ref{rem:reconstruction-stability}).
\end{itemize}
This mechanism is standard in integrable operator theory
(see~\cite{ItsIzergin1990,Kato1966}) and is implemented here
under Assumptions~A and~B.
\end{remark}

\begin{remark}[Comparison to random matrix universality]
\label{rem:rmt-comparison}
Unlike random matrix ensembles, where Airy universality arises
from orthogonal polynomial asymptotics and the $g$-function is
canonically determined by the equilibrium measure, the prolate
operator is deterministic and compact (but not finite rank),
and no orthogonal polynomial structure is available.
The present reduction shows that the same integrable mechanism
--- via the explicit IIKS structure of the sinc kernel --- governs
the transition regime, but the $g$-function must be found as the
solution of an independent nonlinear variational problem rather than
being determined automatically by a known measure.
This combination of IIKS structure with Widom phase asymptotics
constitutes a new class of spectral universality problems, distinct
from both random matrix theory and classical integrable systems.
\end{remark}

%%%=================================================================
\section{Open programme}
\label{sec:open}
%%%=================================================================

\begin{problem}[Assumption~A: Solvability of the nonlinear $g$-problem]
\label{prob:g-existence}
Find an analytic function $g_c:\mathbb{C}\setminus(-\infty,1]\to\mathbb{C}$
with $g_c(z)=\log z + O(z^{-1})$ satisfying the nonlinear matching
condition
\[
  2\operatorname{Re}(g_c^+(t)) = \operatorname{Re}(\phi_c(t)),
  \quad t\in(-1,1),
\]
for all sufficiently large $c$, where $\phi_c = \phi_c[g_c]$
is the Widom edge phase (implicitly depending on $g_c$).
Establish existence (and, if possible, uniqueness) of such a $g_c$.

By Lemma~\ref{lem:semicircle-mismatch}, no fixed candidate from
random matrix theory satisfies this condition for large $c$.
The correct $g_c$ must be genuinely $c$-dependent, tracking the
$O(c)$ growth of the Widom phase.
A natural approach is to linearise the nonlinear coupling
(treating $\phi_c$ as given for a fixed iterate) and apply the
classical theory of singular integral equations on an interval
(Muskhelishvili~\cite{Muskhelishvili1953}), iterating to a fixed
point via the operator $\mathcal{T}_c$ of
Remark~\ref{rem:fixpoint}.
Existence of a fixed point of $\mathcal{T}_c$ is the central
open question.
\end{problem}

\begin{problem}[Assumption~B: Steepest-descent geometry]
\label{prob:rh-programme}
Given a compatible $g_c$ solving Problem~\ref{prob:g-existence},
verify the geometric conditions of Assumption~B:
\begin{enumerate}[label=(\roman*)]
\item \emph{Sign condition:}
  confirm that lens contours $\Sigma_c^\pm$ satisfying
  $\operatorname{Re}(2g_c - \phi_c) < -\kappa < 0$ exist and can be
  chosen to be uniformly Lipschitz in $c$;
\item \emph{Gradient bound:}
  establish $|\nabla_z \operatorname{Re}(2g_c - \phi_c)| \le C$
  on a tubular neighbourhood of $\Sigma_c^\pm \setminus D_\epsilon$,
  uniformly in $c$;
\item \emph{No additional critical points:}
  confirm that $\operatorname{Re}(2g_c - \phi_c)$ has no critical
  points on $\Sigma_c^\pm \setminus D_\epsilon$, so that the $A_2$
  coalescing saddle at $z=0$ is the unique relevant saddle.
\end{enumerate}
Once Problems~\ref{prob:g-existence} and~\ref{prob:rh-programme} are
resolved, Proposition~\ref{prop:small-norm} gives $R_c = I + O(c^{-1/3})$,
Remark~\ref{rem:reconstruction-stability} supplies the Cauchy operator
and resolvent bounds, Proposition~\ref{prop:kernel-from-rhp}
gives $\mathfrak{S}_1$-convergence, and
Conjecture~\ref{conj:airy-transition} becomes a theorem.
\end{problem}

\begin{remark}[Dependency graph]
The logical structure of the paper is:
\[
  \underbrace{\text{IIKS exact}}_{\text{rigorous}}
  \;\Rightarrow\;
  \underbrace{\text{RHP}~(\ref{def:bare-rhp})}_{\text{rigorous}}
  \;\xrightarrow{\text{Ass.~A}}
  \underbrace{J_c^T \text{ oscillatory}}_{\text{conditional}}
  \;\xrightarrow{\text{Ass.~B}}
  \underbrace{R_c\to I}_{\text{conditional}}
  \;\Rightarrow\;
  \underbrace{Y_c\to Y_{\mathrm{Ai}}}_{\text{conditional}}
  \;\xrightarrow{\text{Rem.}~\ref{rem:reconstruction-stability}}
  \underbrace{\norm{\widetilde{K}_c - K_{\mathrm{Ai}}}_{\mathfrak{S}_1}
  \to 0}_{\text{conditional}}
  \;\Rightarrow\;
  \underbrace{\text{Cors.~\ref{cor:gap}--\ref{cor:paper13b}}}_{\text{conditional}}.
\]
Assumption~A is the structural bottleneck: it controls whether the
Deift--Zhou programme can be initiated at all, and its nonlinear
character means it has no counterpart in the classical OP setting.
Assumption~B is the analytic completion: standard in form, but
dependent on the unknown solution of Assumption~A.
The completeness of the reduction --- that Assumptions~A and~B
together suffice --- is the content of
Proposition~\ref{prop:reduction-completeness}.
\end{remark}

%%%=================================================================
\section{Conclusion and outlook}
\label{sec:conclusion}
%%%=================================================================

The present work establishes a complete Riemann--Hilbert reduction
of the prolate Airy transition problem to two logically separated
analytic inputs: a structural nonlinear variational problem
(Assumption~A) and a geometric steepest-descent condition
(Assumption~B).

This reduction has four distinguishing features:
\begin{enumerate}[label=(\roman*)]
\item \emph{Exact integrable structure.}
The sinc kernel admits an explicit rank-2 IIKS representation,
providing a natural Riemann--Hilbert formulation from the outset.
This is not an approximation but an algebraic identity.

\item \emph{Precise asymptotic control.}
The effective edge phase is analysed to all orders in $c^{-1/3}$,
revealing a coalescing $A_2$ saddle structure.
The cubic Airy term arises only after integrating the Widom phase
(Lemma~\ref{lem:edge-normal-form}), not from the Fourier
representation of the kernel itself --- a key structural distinction
that explains why Airy universality is a global RHP phenomenon and
cannot be seen locally.

\item \emph{Variational identification of the core difficulty.}
The central open problem is identified as the solvability of a
nonlinear scalar Riemann--Hilbert matching condition for the
compatible $g$-function (Assumption~A,
Problem~\ref{prob:g-existence}).
By Lemma~\ref{lem:semicircle-mismatch}, no fixed candidate from
random matrix theory can serve as a solution; the correct $g_c$
is genuinely $c$-dependent and must track the $O(c)$ growth of
the Widom phase.
This is the structural difficulty that has no analogue in the
classical OP setting, where $g$ is determined canonically.

\item \emph{Explicit and modular programme.}
All components of the Deift--Zhou steepest descent analysis
are constructed explicitly and the two open inputs are
logically separated: Assumption~A (structural, nonlinear) must
be established first; Assumption~B (geometric, conditional on~A)
is the expected next step.
\end{enumerate}

The main result can be stated concisely:

\begin{center}
\fbox{\parbox{0.9\linewidth}{\centering\itshape
Airy universality for the prolate concentration operator in the
transition regime reduces to, and is implied by, the existence of
a $g$-function $g_c$ satisfying the nonlinear matching condition
$2\operatorname{Re}(g_c^+(t)) = \operatorname{Re}(\phi_c(t))$
on $(-1,1)$ (Assumption~A), together with the steepest-descent
geometry of Assumption~B.
The entire analytic difficulty of the Airy transition is encoded
in the solvability of this nonlinear scalar Riemann--Hilbert
equation, whose solvability is currently open.}}
\end{center}

Under Assumptions~A and~B, the paper establishes:
\begin{itemize}
\item Fredholm determinant universality (Theorem~\ref{thm:universality}),
\item $c^{-1/3}$ gap asymptotics (Corollary~\ref{cor:gap}),
\item local spectral statistics and stable quasimodes
  (Corollaries~\ref{cor:counting}--\ref{cor:quasimodes}),
\item conditional completion of the sharp transition analysis of
  Paper~XIII (Corollary~\ref{cor:paper13b}).
\end{itemize}

The combination of IIKS structure with Widom phase asymptotics
identifies a new class of spectral universality problems: not
random matrix ensembles, not classical integrable systems, but
deterministic compact operators whose Airy transition is governed
by a nonlinear compatibility equation between two a priori
unrelated analytic objects --- the equilibrium potential and the
Widom edge phase.
The variational $g$-problem (Problem~\ref{prob:g-existence}) is
the natural and necessary next step toward a fully rigorous proof
of Airy universality for prolate spheroidal wave functions.

More broadly, the present framework suggests a general principle:
for deterministic compact operators arising from time--band
limiting, edge universality at the Airy level is governed not
by orthogonal polynomial machinery but by a self-consistent
$g$-function problem of the type identified here.
This may be the correct framework for a general theory of
spectral universality beyond random matrix ensembles.

%%%=================================================================
\begin{thebibliography}{99}

\bibitem{Slepian1961}
D.~Slepian, H.~O.~Pollak,
\textit{Prolate spheroidal wave functions, Fourier analysis and
uncertainty~I},
Bell System Tech. J. \textbf{40} (1961), 43--63.

\bibitem{Widom1964}
H.~Widom,
\textit{Asymptotic behavior of the eigenvalues of certain integral
equations},
Trans. Amer. Math. Soc. \textbf{109} (1964), 278--295.

\bibitem{Deift1999}
P.~Deift,
\textit{Orthogonal Polynomials and Random Matrices: A Riemann--Hilbert
Approach},
Courant Lecture Notes, AMS, 1999.

\bibitem{TraWid1994}
C.~A.~Tracy, H.~Widom,
\textit{Level-spacing distributions and the Airy kernel},
Comm. Math. Phys. \textbf{159} (1994), 151--174.

\bibitem{ItsIzergin1990}
A.~R.~Its, A.~G.~Izergin, V.~E.~Korepin, N.~A.~Slavnov,
\textit{Differential equations for quantum correlation functions},
Int. J. Mod. Phys. B \textbf{4} (1990), 1003--1037.

\bibitem{Osipov2013}
A.~Osipov, V.~Rokhlin, H.~Xiao,
\textit{Prolate Spheroidal Wave Functions of Order Zero},
Springer, 2013.

\bibitem{Kato1966}
T.~Kato,
\textit{Perturbation Theory for Linear Operators},
Springer, 1966.

\bibitem{Muskhelishvili1953}
N.~I.~Muskhelishvili,
\textit{Singular Integral Equations},
Noordhoff, Groningen, 1953.

\bibitem{PaperVIII}
[Author], \textit{Paper~VIII}, 2026.

\bibitem{PaperXII}
[Author], \textit{Paper~XII}, 2026.

\bibitem{PaperXIII}
[Author], \textit{Paper~XIII (paper13b)}, 2026.

\end{thebibliography}

\end{document}
